\documentclass[12pt]{article}
\usepackage[utf8]{inputenc}
\usepackage[T1]{fontenc}
\usepackage{hyperref}
\title{Vale Language Resources for K-12 Curriculum Development}
\author{Compiled by Claude Sonnet 5 (claude-sonnet-5) for Andr\'as Kornai}
\date{}
\begin{document}
\maketitle
\section{Introduction}
Vale (ISO 639-3: vae) is a Central Sudanic language of the Central African Republic with about 5,400 speakers (a 1996 figure). Search found almost nothing specific to Vale; honest reporting means this dossier is very short.
\section{Primary Educational Resources}
None found.
\section{Linguistic and Reference Materials}
\begin{itemize}
\item Vale language (Wikipedia): classification in the Ruto--Vale group and basic facts. \url{https://en.wikipedia.org/wiki/Vale_language}
\end{itemize}
No Vale dictionary or grammar was found. A Ngbaka--French dictionary appeared in results but concerns a different, Ubangian language and is not a Vale resource.
\section{Religious and Community Resources}
None found.
\section{Digital and Audio Resources}
\begin{itemize}
\item Global Recordings Network page for Vale (audio listing). \url{https://globalrecordings.net/en/language/vae}
\item Joshua Project people-group profile, Vale in Central African Republic (demographic information). \url{https://joshuaproject.net/index.php/people_groups/15764}
\end{itemize}
\section{Curriculum Development Context}
No evidence of a written standard or school use was found. Any curriculum would require new documentation work.
\end{document}